\documentclass[11pt]{article}
\usepackage[T1]{fontenc}
\usepackage[utf8]{inputenc}
\usepackage{lmodern,microtype,amsmath,amssymb,amsthm,mathtools}
\usepackage[margin=1in]{geometry}
\usepackage{booktabs,longtable,array,xcolor,enumitem,fancyhdr,listings}
\usepackage{xurl}
\usepackage[hypertexnames=false,colorlinks=true,linkcolor=blue!45!black,urlcolor=blue!45!black,citecolor=blue!45!black]{hyperref}
\usepackage{bookmark}
\hypersetup{pdftitle={Four-rule iPRA with reusable clauses},
 pdfauthor={Prepared for Vladimir Reshetnikov},
 pdfsubject={Explicit axiomatization, bi-interpretation, and a complete proof checker}}
\definecolor{ink}{HTML}{173A52}
\definecolor{muted}{HTML}{526573}
\setlength{\parindent}{0pt}
\setlength{\parskip}{0.55em}
\setlength{\emergencystretch}{2em}
\setlist{nosep,leftmargin=1.7em}
\pagestyle{fancy}
\fancyhf{}
\fancyhead[L]{\small\textcolor{muted}{Four-rule iPRA}}
\fancyhead[R]{\small\textcolor{muted}{Clause contexts}}
\fancyfoot[C]{\small\thepage}
\setlength{\headheight}{14pt}
\newtheorem{theorem}{Theorem}[section]
\newtheorem{lemma}[theorem]{Lemma}
\newtheorem{proposition}[theorem]{Proposition}
\newtheorem{corollary}[theorem]{Corollary}
\theoremstyle{definition}
\newtheorem{definition}[theorem]{Definition}
\newtheorem{remark}[theorem]{Remark}
\newcommand{\RC}{\mathcal C_4}
\newcommand{\Old}{\mathcal C_5}
\newcommand{\ipra}{\mathsf{iPRA}}
\newcommand{\RR}{\mathcal R}
\newcommand{\nf}{\mathsf N}
\newcommand{\cl}[1]{\langle #1\rangle}
\newcommand{\FV}{\operatorname{FV}}
\lstset{basicstyle=\ttfamily\fontsize{8}{9}\selectfont,columns=fullflexible,keepspaces=true,
  breaklines=true,breakatwhitespace=false,showstringspaces=false,
  frame=single,rulecolor=\color{black!15},xleftmargin=3pt,xrightmargin=3pt}
\newcommand{\SelectedBytes}{2282}
\newcommand{\SelectedTokens}{666}
\newcommand{\SelectedAST}{770}
\begin{document}
\thispagestyle{empty}
{\Large\bfseries\color{ink} Four-rule iPRA with reusable clauses}\par
\vspace{0.4em}
{\large An explicit presentation, exact proof translations, and a smaller complete checker}\par
\vspace{0.8em}
Prepared for Vladimir Reshetnikov\hfill 11 November 2026\par
\textcolor{muted}{Comparison pinned to ProveIt commit \texttt{443cff4956}.}\par
\vspace{0.7em}\hrule\vspace{0.8em}

\begin{abstract}
We give an incremental simplification of the current ConciseIPRA presentation.
One fixed reusable clause, $\RR=[\,\,;y.(y=y)]$, and four inference schemes
--- Use, Step, Open, and Close --- present a theory bi-interpretable and
equiconsistent with full first-order intuitionistic primitive recursive
arithmetic. Universal reflexivity, ordinary hypotheses, and conditional
closures all use the same clause mechanism. Numerical terms are zero-normal
scalar iterators; all normalization is included in the checker.

We prove exact translations to and from the repository's five-pattern
zero-normal clause calculus, recover its arithmetic and logical interface,
and make both interpretation composites explicit. A complete standalone
Python checker, bidirectional certificate translators, negative controls,
and reproducible source measurements accompany the paper. The selected
checker has \SelectedBytes\ normalized bytes, \SelectedTokens\ significant
tokens, and \SelectedAST\ AST nodes, compared with $2950/829/902$ for the
pinned transparent five-pattern checker. The comparison includes parsing,
syntax and scope validation, hereditary substitution, all rules, and the
acceptance interface.

The improvement is a presentation and implementation reduction. Reflexivity
remains a counted axiom clause. No new minimum over the repository's compact
codecs, no absolute minimality theorem, and no new Lean verification are claimed.
\end{abstract}

\section{Result and comparison boundary}
\label{sec:main-result}

The target is \emph{full first-order} $\ipra$: intuitionistic first-order
logic with equality, a symbol for each fixed primitive-recursive definition,
the defining equations, successor separation and injectivity, and induction
for quantifier-free formulas. In the full primitive-recursive language,
equation induction and quantifier-free induction are interderivable via
characteristic functions. The first-order formulation matters: agreement on
equations or on provably total functions alone would not establish the requested
bi-interpretation. See also Salehi~\cite{salehi} for the first-order equation-
induction convention.

The pinned repository already contains scalar iterators, the paired binder
$D$, clause contexts, a single successor Step that recovers equality, and
zero-normal syntax~\cite{baseline,uniformstep,nativezerostep}. These are
inherited ingredients. The new interface shares clause application between
three previously different uses, and the implementation removes repeated
proof-tree annotations through a goal-directed packet format.

\begin{theorem}[Main result]
\label{thm:main}
The calculus $\RC$ specified in Sections~\ref{sec:syntax}--\ref{sec:rules}
has the following properties.
\begin{enumerate}
\item Its public consequences are exactly those of the pinned five-pattern
zero-normal clause calculus $\Old$, under the identity translation on
numerical terms and public formulas.
\item Its ordinary realization is bi-interpretable with $\ipra$. Both
interpretation domains are the entire numerical domain, equality is actual
equality, and the two numerical composite isomorphisms have graph $x=y$.
Both formula composites are provably equivalent to the original formula.
\item $\RC$ proves its designated contradiction $0=S0$ if and only if
$\ipra$ proves falsity. Finite proof translations are effective.
\item The delivered standalone checker accepts exactly its closed theorems
when run with sufficient resources, with the documented certificate syntax.
\end{enumerate}
\end{theorem}

Here ``exactly'' describes the mathematical checker algorithm on finite
certificates. Ordinary Python recursion and memory limits can reject a very
large valid certificate. They do not create accepted false conclusions.
Neither theory is asserted to prove its own consistency.

\paragraph{How much has actually been simplified?}
Four inference schemes and one initial axiom clause replace five inference
schemes, one of which was reflexivity. Thus counting ``four'' alone would
overstate the result. The substantive sharing is that a single operation
handles reflexivity, hypotheses, and conditional clauses. The complete
checker comparison, including the initial clause's compiled behavior, is
the evidence for a size reduction. A stronger matched control and every
selected representation cost are reported in Section~\ref{sec:measurements}.

The repository also records a $516$-byte compact profile and a $788$-byte
public-goal validated profile~\cite{baseline,compact}. These have different,
heavily fused certificate encodings. The present programs do not improve
those numerical records. There is no proved absolute lower bound for all
axiomatizations or implementations in this report.

\section{Numerical syntax and public formulas}
\label{sec:syntax}

\subsection{Zero-normal scalar iteration}

There is one numerical sort. Start with raw terms
\[
 t ::= x\mid 0\mid St\mid I(n,a;u.h).
\]
Only $u$ is bound in $h$; it is not bound in $n$ or $a$.
Other parameters in the update remain fixed. Define the smart constructor
\[
 \widehat I(n,a;u.h)=
 \begin{cases}a,& n\text{ is the literal }0,\\
 I(n,a;u.h),&\text{otherwise}.
 \end{cases}
\]
Define $\nf$ recursively on the syntax: it fixes variables and zero,
commutes with $S$, and rebuilds each iterator with $\widehat I$ after
normalizing its three children. The actual terms of $\RC$ are those fixed
by $\nf$. There is no iterator whose counter is literally zero anywhere
in such a term. In a rule, substitution means capture-free substitution
followed by $\nf$.

This builds in only $I(0,a;u.h)=a$. It does \emph{not} normalize successor
iterations, execute an arbitrary primitive-recursive function, or appeal to
a numerical evaluator. The counter may be any term, including an iterator.
The operation terminates by structural recursion on the finite input tree:
the zero contraction returns its already-normalized seed and does not
duplicate or reprocess an update.

\begin{lemma}[Zero-normal substitution]
\label{lem:nf}
Normalization is idempotent, does not introduce free variables, and satisfies
\[
 \nf(t[\sigma])=\nf\bigl(\nf(t)[\nf\circ\sigma]\bigr).
\]
Consequently hereditary substitution is coherent under composition.
\end{lemma}
\begin{proof}
Induct on the raw term. Variables, zero, and successors are immediate.
For an iterator, normalize the counter, seed, and update first. The two
sides have equal normalized children by induction. If their common counter
is literal zero both return the common normalized seed; otherwise both
return the same iterator. Under its update, lift the substitution past its
one bound variable. Idempotence is the same structural argument with the
identity substitution. No semantic equality test is used.
\end{proof}

\subsection{The Eq/\texorpdfstring{$D$}{D} language}

Public formulas are
\[
 A ::= (s=t)\mid D_{x,y}A.
\]
The variables $x,y$ bound by $D$ are distinct and are renamed before
substitution. The intended ordinary reading is
\begin{equation}
 (D_{x,y}B)^\circ
   =\exists x\bigl(B^\circ(x,0)\to\forall y\,B^\circ(x,y)\bigr).
 \label{eq:Dreading}
\end{equation}
Equality is actual numerical equality. The designated contradiction is
$\bot_* := (0=S0)$. There are no primitive symbols for ordinary implication,
conjunction, disjunction, universal quantification, or falsity in this native
grammar. Their reconstruction is given in Section~\ref{sec:bridge} and the
inherited logical appendix.

\paragraph{A useful exact representation.}
The grammar is a unary spine: every public formula is uniquely
$D^k(s=t)$ after fixing a binding convention. Thus it can be represented
as a ranked equation $(k,s,t)$. With de Bruijn indices, the endpoints have
$2k$ bound numerical slots. This is a bijective encoding, not a prenex
normal-form theorem and not a claim that arbitrary ordinary formulas have
a single quantifier prefix. Each $D$ has the conditional reading
\eqref{eq:Dreading}. We implemented and tested this representation, but
an equally careful recursive Eq/$D$ implementation was smaller. It is
therefore reported as a comparison, not used to manufacture a size claim.

\section{One fixed clause and four inference schemes}
\label{sec:rules}

A reusable clause has zero or one antecedent:
\[
 [A_1,\ldots,A_m;y.B(y)],\qquad m\in\{0,1\},\quad
 y\notin\bigcup_j\FV(A_j).
\]
Its ordinary reading is $\forall y B^\circ(y)$ for $m=0$ and
$A_1^\circ\to\forall y B^\circ(y)$ for $m=1$. Contexts are finite lists
of reusable clauses. Every context has access to the distinguished clause
\begin{equation}
                 \RR=[\,\,;y.(y=y)].
                 \tag{R}\label{eq:R}
\end{equation}
An ordinary hypothesis $A$ is the constant unconditional clause
$\cl{A}:=[\,\,;y.A]$, with $y$ absent from $A$. Lists, rather than sets,
make context occurrence indices unambiguous in certificates.

The complete rule inventory is the following. All displayed substitutions
include the zero normalization of Section~\ref{sec:syntax}.

\paragraph{Use.}
For a retained clause and a well-scoped numerical term $t$,
\[
 \frac{[A_1,\ldots,A_m;y.B]\in\Gamma\qquad
       \Gamma\vdash A_1\quad\cdots\quad\Gamma\vdash A_m}
      {\Gamma\vdash B(t)}.
 \tag{Use}
\]
The premise list is empty for $m=0$. In particular Use of $\RR$ at $t$
proves $t=t$.

\paragraph{Step.}
\[
 \frac{\Gamma\vdash\widehat I(c,a;u.h)=t\qquad
       \Gamma\vdash I(Sc,a;u.h)=r}
      {\Gamma\vdash r=h[t/u]}.
 \tag{Step}\label{eq:Step}
\]
The counter, seed, and update agree literally in the two source iterators,
up to the built-in zero contraction of the first one. The update is a fixed
term template, not a quantified function variable.

\paragraph{Open.}
\[
 \frac{\Gamma\vdash D_{x,y}B(x,y)\qquad
       \Gamma,[B(b,0);y.B(b,y)]\vdash A}
      {\Gamma\vdash A}.
 \tag{Open}
\]
The numerical eigenparameter $b$ is fresh for $\Gamma$, $A$, and all
fixed parameters of the displayed matrix. The added clause is discharged.

\paragraph{Close.}
\[
 \frac{\Gamma,\cl{B(t,\eta)}\vdash B(t,Sz)}
      {\Gamma\vdash D_{x,y}B(x,y)},\qquad
 \eta=
 \begin{cases}
 z,& B\text{ is a literal equation and }t\text{ is literal }0,\\
 0,&\text{otherwise}.
 \end{cases}
 \tag{Close}\label{eq:Close}
\]
Here $z$ is fresh for $\Gamma$, $t$, and the fixed parameters of $B$.
The local unconditional clause is discharged. ``Literal equation'' excludes
an equation beneath a $D$ binder. The witness test is a syntactic test,
not a test of provable equality with zero.

\paragraph{Nothing else is primitive.}
There is no other axiom, equality rule, induction rule, logical rule,
structural rule, or computation equation. Weakening, renaming, and cut are
admissible; they can be proved by translation or the usual induction on
derivations, with local eigenparameters renamed. General numerical
substitution is admissible via the equivalent unrestricted basis, but
should not be justified by blindly copying Close: a substitution can turn
a nonzero-looking witness into literal zero and change its selected mode.

\section{Recovery of computation, equality, and both introductions}
\label{sec:recovery}

\subsection{One successor Step supplies the arithmetic interface}

For brevity, write $h(t)$ for normalized substitution. All terms below are
well-scoped and zero-normal.

\begin{lemma}[Computation and the equality pivot]
\label{lem:pivot}
The calculus derives the successor equation
\[
 I(Sc,a;u.h)=h(\widehat I(c,a;u.h))
\]
and the pivot inference
\[
             s=t,\ s=r\ \vdash\ r=t.
\]
Neither derivation uses induction or a logical rule.
\end{lemma}
\begin{proof}
Successor computation is Step on the two reflexivities supplied by Use of
$\RR$. For the pivot let $J=I(S0,s;u.u)$. Step on $s=r$ and $J=J$ gives
$J=r$. A second Step on $s=t$ and $J=r$ gives $r=t$.
\end{proof}

Symmetry is the pivot on $s=s$ and $s=t$. Given $s=t$ and $t=u$, reverse
the first equality and apply the pivot to obtain $s=u$.
For transport in one term, assume $s=t$ and $h(s)=v$. Put
$K=I(S0,s;u.h)$. Step gives $K=h(s)$ from reflexivities. Transitivity gives
$K=v$. Step on $s=t$ and $K=v$ gives $v=h(t)$. Two such term transports
give transport in an equation template. Transport in compound formulas
follows by induction on the $D$ spine using Open and case introduction.

The zero equation is definitional, with its full normalizer counted in the
implementation. The successor equation is derived. Universal reflexivity
has not disappeared: its source is exactly clause~\eqref{eq:R}.

\subsection{Why the single Close rule is sufficient}

The familiar case introduction would allow
$\Gamma,B(t,0)\vdash B(t,Sz)$ for every $t$; equation progress would allow
$\Gamma,E(t,z)\vdash E(t,Sz)$. Close selects one mode deterministically.
The apparent missing atomic case at literal witness zero is recoverable.

Let
\[
                       q=I(S0,0;u.u).
\]
It is not syntactically zero, and Lemma~\ref{lem:pivot} gives $q=0$.
Suppose $\Gamma,E(0,0)\vdash E(0,Sz)$. First use Close on the matrix
$E(0,y)$ with an unused witness coordinate and witness $q$. This uses
case mode, and gives $D_{v,y}E(0,y)$. Open it to obtain the clause
$[E(0,0);y.E(0,y)]$. Now introduce $D_{x,y}E(x,y)$ with witness $q$:
its assumption $E(q,0)$ transports to $E(0,0)$; Use supplies $E(0,Sz)$;
transport back gives $E(q,Sz)$. Close and then Open discharge the added
data. The desired atomic case conclusion follows.

Equation progress for a fixed matrix $E(t,y)$ is introduced first with an
unused witness coordinate and literal witness zero. Its Open clause gives
the conditional closure of $E(t,y)$. This supplies every witnessed progress
instance through the now unrestricted case rule. In particular, from
\[
 \Gamma\vdash E(0),\qquad \Gamma,E(z)\vdash E(Sz)
\]
one derives $\Gamma\vdash E(t)$ for any $t$ by Close, Open, and Use.
Numerical cases for arbitrary formulas are obtained in the same way from
their zero proof and fresh successor proof using case mode.

These constructions explain why Close is sufficient and why extending
progress to arbitrary compound matrices would be a different theory.
The latter would give induction for arbitrary formulas by exactly the
same Open--Use argument.

\section{Exact proof equivalence with the pinned calculus}
\label{sec:equivalence}

The matched source $\Old$ uses the same zero-normal terms, public Eq/$D$
formulas, Step, Open, and witness-selected Close. It has explicit reflexivity
and conditional clause application. A local ordinary hypothesis $A$ is
stored as $[0=0;y.A]$ and needs a proof of its neutral antecedent.
This is the complete checker in
\nolinkurl{Research/uniform_step/code/checker.py} at the pinned commit.

\begin{theorem}[Clause factorization]
\label{thm:factorization}
There are effective translations of finite proofs in both directions
between $\Old$ and $\RC$, preserving each public conclusion and retained
side assumptions under the clause translations below. In particular,
their closed theorem sets are identical.
\end{theorem}
\begin{proof}
For the forward translation, insert $\RR$ as the first new context
occurrence. Replace source reflexivity at $t$ by Use of that occurrence
at $t$ with no premises. Mark each source context occurrence according to
whether Open or Close introduced it. Open occurrences remain conditional.
A Close occurrence $[0=0;y.A]$ becomes $\cl{A}$; a Use of that occurrence
discards the source proof of $0=0$. Translate other premises recursively.
Step, Open, and Close retain the same public formulas and witnesses.

For the reverse translation, a Use of the distinguished $\RR$ becomes
source reflexivity. Every other unconditional clause $[\,\,;y.B]$ is
represented by $[0=0;y.B]$, and its application gets the source proof
$0=0$ by reflexivity. Conditional clauses and their applications are
unchanged. Again translate the other rules recursively.

The induction tracks context \emph{occurrences}, not only their formula
contents. Therefore equal-looking clauses do not confuse the conversion,
and discharge and eigenparameter conditions are preserved. No equality
proof, computation, or logical rule is imported as an oracle. The proof
does not need substitution stability of the witness-selected rule.
\end{proof}

\paragraph{Representation changes.}
Removing the numerical constructor tags is a bijection: $St$ becomes
$[t]$ and $I(n,a;u.h)$ becomes $[n,a,h]$. Public equality becomes the
endpoint pair $[s,t]$; the paired formula remains $[\texttt{d},B]$.
The types make these shapes disjoint: a numerical term is an integer or
a list of length one or three, and a formula is a list of length two.
The tag \texttt{d} is never a numerical term. Hence one opening routine
can traverse all these objects with the correct binder increments.
Structural induction proves that the encoding commutes with normalized
opening, so it transports the entire proof system without semantic change.

\section{Full intuitionistic logic and the bi-interpretation}
\label{sec:bridge}

\subsection{The ordinary realization}

Let $T_I$ be intuitionistic first-order logic with equality over the staged
scalar iterator language. For each fixed earlier-stage update $h(u,\vec p)$,
introduce a symbol $i_h(n,a,\vec p)$ with
\[
 i_h(0,a,\vec p)=a,\qquad
 i_h(Sn,a,\vec p)=h(i_h(n,a,\vec p),\vec p).
\]
Add numerical coverage $\forall x(x=0\vee\exists y\,x=Sy)$, the separation
bridge $(0=S0)\to\bot$, and contextual induction for literal equations.
This uses fixed finite function descriptions, with no uniform evaluator
over function codes.

Interpret $\RC$ by actual equality, \eqref{eq:Dreading}, and the clause
readings in Section~\ref{sec:rules}. Use and Open are ordinary logical
inferences. Step follows from the defining equation and equality. Case
Close is sound by coverage; its progress case is sound by equation
induction under a temporarily assumed base. The initial clause is equality
reflexivity. This proves soundness for the \emph{whole} ordinary domain.

The explicit separation bridge is essential in the ordinary theory:
positive iterator equations alone admit a one-element structure. Native
falsity is $0=S0$; ordinary falsity is $\bot$. Equiconsistency must preserve
this distinction.

\subsection{Logical macros and their composites}

Set $C(v,a,b)=I(v,a;u.b)$, with $u$ unused, so $C(0,a,b)=a$ and
$C(Sv,a,b)=b$. The gate $G_v A$ replaces a terminal equation $s=t$ by
$C(v,0,s)=C(v,0,t)$ and commutes with $D$. Thus $G_0A$ is derivable and
$G_{Sv}A$ is interderivable with $A$. Padding
$A\dashv\vdash D_{u,y}G_yA$ aligns formula depths.

Using this padding, define a selector $M_v(A,B)$ with
$M_0(A,B)\dashv\vdash A$ and $M_{Sv}(A,B)\dashv\vdash B$:
at equations use $C$ on both endpoints, and at equal $D$ depths recurse
under their aligned binders. When one operand is an equation and the other
has an outer $D$, pad the equation with $G_y$ and recurse under that $D$.
The maximum input depth decreases at each call. Let
$p(z)=I(z,0;u.C(u,1,0))$; fixed computations give $p(1)=1$, $p(2)=0$.
Then
\[
\begin{aligned}
 A\Rightarrow_*B&=D_{u,z}M_z(A,B),\\
 A\wedge_*B&=D_{u,z}G_zM_{p(z)}(A,B),\\
 \exists_*x A(x)&=D_{x,y}G_yA(x),\\
 A\vee_*B&=\exists_*x M_x(A,B),\\
 \forall_*x A(x)&=A(0)\wedge_*D_{u,z}G_zA(z),\\
 \bot_*&=(0=1).
\end{aligned}
\]
All added bound variables are fresh, and $u$ is unused. The inherited
logical appendix supplies the introduction and elimination derivations.
For example, implication's zero instance supplies its assumption and its
successor instances supply the consequent. Existential introduction uses
the chosen witness and the gate; Open, the true zero gate, and Use at $1$
give existential elimination. Numerical cases supply disjunction
elimination and the missing zero case of universal elimination.

For the crucial native formula composite one proves
\[
 D_{x,y}B\dashv\vdash
 \exists_*x\bigl(B(x,0)\Rightarrow_*\forall_*y B(x,y)\bigr).
\]
Forward, Open supplies its conditional clause; introduce the three logical
macros around that clause. Backward, eliminate the existential at a fresh
witness and apply case Close using implication and universal elimination.
Induction on finite formulas proves both logical composites. This is more
than numerical conservativity.

\subsection{Arithmetic translations and identity composites}

In $\ipra$, interpret a scalar iterator by primitive recursion whose update
ignores the recursion counter. Conversely, the inherited scalar development
establishes, in this order: a predecessor by equation induction and cases;
truncated subtraction and a characteristic of actual equality; quantifier-
free induction; and pairing with projections. The detailed proofs are
reproduced in the arithmetic appendices~\cite{erasure,forward,progress}.
This order avoids assuming pairing or equality reflection in the predecessor
bootstrap.

For a primitive recursion
\[
 f(0,\vec p)=g(\vec p),\qquad
 f(Sn,\vec p)=h(n,f(n,\vec p),\vec p),
\]
let $\pi,L,H$ be the derived pairing and projections. Its scalar translation is
\[
\begin{aligned}
 V(n,\vec p)&=I\bigl(n,\pi(0,g^*(\vec p));
 u.\pi(SL(u),h^*(L(u),H(u),\vec p))\bigr),\\
 f^*(n,\vec p)&=H(V(n,\vec p)).
\end{aligned}
\]
Equation induction proves $L(V(n,\vec p))=n$; the projection equations then
prove both defining equations for $f^*$. Initial functions and composition
are direct. This is external recursion on each finite definition.

Two scalar iterators with equal seeds and provably equal updates are equal
at every internal counter by equation induction. The analogous uniqueness
proof works for ordinary primitive recursion. External induction on a
finite function description therefore proves that its double translation
equals the original operation internally. The interpretations use the
whole domain and actual equality; $x=y$ is consequently a provably total
bijective graph preserving all operations in both composites. Combine this
with the preceding formula composites and the falsity bridge to prove
Theorem~\ref{thm:main}(2)--(3).

The repository formalizes the inherited bridge in
\nolinkurl{BiInterpretation.lean} and \nolinkurl{NativeBiInterpretation.lean};
its zero-normal Step interface is \nolinkurl{NativeZeroStep.lean}~\cite{biinterp,
nativebi,nativezerostep}. We inspected the pinned sources. We did not rebuild
that Lean project in this run. The new clause factorization and packet
checker correctness here are written proofs with executable cross-checks,
not additional Lean theorems.

\section{The complete checker and its certificates}
\label{sec:checker}

\subsection{Binding and validation}

The selected checker uses exact integers and lists. Zero is integer $0$;
positive integers are de Bruijn indices; negative integers are fresh free
levels. At numerical scope $q$ and bound depth $d$, an integer is valid
exactly when $-q\le a\le d$. Thus Python booleans and floats are rejected
even though some of them compare equal to integers. Successor is $[t]$;
an iterator is $[n,a,h]$, with $n\ne0$ syntactically and the body checked
at depth $d+1$. A formula is $[s,t]$ or $[\texttt{d},B]$; its $D$ matrix
is checked at depth $d+2$.

In a paired matrix, index $1$ denotes $y$ and index $2$ denotes $x$,
before accounting for any further inner binders. At each reachable packet
checker state, the numerical scope equals the length of the explicit
context: the fixed $\RR$ is ambient, and each Open or Close adds exactly
one clause and one new numerical level. This invariant removes a separate
scope argument without identifying arbitrary contexts with variable scopes.

Opening at cutoff $d$ replaces index $d+1$ by the supplied term, decreases
larger indices by one, and leaves smaller indices and negative levels
unchanged. It increments the cutoff by one under an iterator body and by
two under $D$. Rebuilding an iterator with counter zero returns its seed.
Inserted terms have no free positive indices, so they require no additional
lifting; their own internal bound indices remain protected. This is a
locally nameless representation with stable negative free levels.

Every intermediate goal is validated before a rule is processed.
For Open, its conclusion is checked before extending the scope; it cannot
mention the new eigenparameter. Close validates its witness in the old
scope. Context clauses are generated only by those rules, so their validity
is maintained internally rather than trusted from input.

\subsection{Goal-directed packets}

The selected complete certificate is a JSON array of two-field packets
followed by a closed goal. Packets are read from right to left. A packet
$[b,t]$ selects its inference from the shape of $b$:
\begin{center}
\begin{tabular}{@{}p{0.24\linewidth}p{0.64\linewidth}@{}}
\toprule
First field $b$ & Meaning of the packet\\
\midrule
Nonnegative integer & Use context clause $b$ at numerical target $t$.\\
$-1$ & Close the current paired goal with witness $t$.\\
A paired formula & Prove $b$, Open it, then prove the retained goal;
the unused second field must be the exact integer $0$.\\
An equation & Prove $b$ and the equation $t$, then check that they are
the two premises of Step for the current goal.\\
\bottomrule
\end{tabular}
\end{center}

A reflexive equation closes immediately by Use of $\RR$, without consuming
a packet. The fixed clause is therefore compiled into a reflexivity test;
its behavior remains counted in the checker. Context indices enumerate
the local clauses, since $\RR$ needs no index in this protocol. All consumed
instructions are checked, and any unconsumed prefix causes rejection.
The companion tree checker implements the explicit four-scheme formulation
with a stored initial clause and a finite list of premises for Use.

\paragraph{Neutral-antecedent lowering.}
An unconditional local clause can be stored internally as a clause with
antecedent $0=0$: the check of that antecedent closes automatically without
consuming input. This is a compilation of empty premise lists, not an
additional axiom and not an omitted proof obligation. It permits the same
premise-checking path for every local clause. The comparison includes this
implementation choice.

\begin{theorem}[Packet checker correctness]
\label{thm:packet}
Successful checking of a goal under a well-scoped context yields a finite
$\RC$ derivation of that goal. Every finite $\RC$ derivation of a closed
theorem, with no ordinary assumptions apart from the ambient $\RR$, has an
accepted packet certificate with the same conclusion.
\end{theorem}
\begin{proof}
For soundness, induct on the finite successful call tree. Automatic closure
is Use of $\RR$. Otherwise one packet is removed; its disjoint shape chooses
one of the four rules. Recursive calls establish its premises, and the
explicit comparisons establish its conclusion and all side conditions.
The pre-extension validation of an Open conclusion gives the eigenvariable
condition. Clause instances use the proved hereditary-opening operation.
Neutral-antecedent lowering uses the fixed reflexivity proof of $0=0$.
The other validation invariants follow by induction on the context.

For completeness, induct on a derivation. Replace any proof of a reflexive
goal by Use of $\RR$. At a nonreflexive goal, emit the packet for its last
rule and serialize its premise proofs in the order in which the checker
will request them. Reverse this instruction sequence in the external array
and append the goal. The suffix-consumption invariant says that one call
consumes exactly its own serialized proof segment and leaves the preceding
segments intact. Prove it by the same induction. Empty-premise Use consumes
no premise segment; its lowered neutral proof also consumes none. Every
rule is represented, and at the root the remaining array is empty.
\end{proof}

This proves correctness of the algorithm in the paper's metatheory.
Finite tests support the implementation correspondence but cannot replace
this universal proof.

\section{Measurements, validation, and limits}
\label{sec:measurements}

All source measurements use CPython 3.12.14, remove documentation strings,
apply \texttt{ast.unparse}, and append one final LF. Significant tokens
exclude comments, layout, encoding, and end markers. AST counts include
the complete parsed acceptance module. Gzip uses level nine and timestamp
zero; it is informational. The exact measured sources and gzip files are
included, together with their hashes and the measurement script.

\begin{center}
\small
\begin{tabular}{@{}p{0.47\linewidth}rrrr@{}}
\toprule
Complete checker & Bytes & Tokens & AST & Gzip 9\\
\midrule
Pinned five-pattern tagged checker & 2950 & 829 & 902 & 1005 \\
Five-rule arity and scope control & 2730 & 753 & 828 & 939 \\
Four-rule explicit tree checker & 2677 & 751 & 823 & 947 \\
Four-rule ranked tree checker & 2692 & 781 & 853 & 964 \\
Four-rule empty-antecedent packets & 2308 & 665 & 770 & 778 \\
Selected four-pattern packet checker & 2282 & 666 & 770 & 774 \\
\bottomrule
\end{tabular}
\end{center}

The five-pattern arity/scope control is deliberately stronger than the
unaltered tagged source: it receives the same routine codec and scope
improvements. In that comparison, clause factorization alone saves only
$53$ bytes, $2$ significant tokens, and $5$ AST nodes; its gzip size grows
by $8$ bytes. The larger reduction comes from the full certificate protocol
and shared validation, not from pretending that a relocated axiom vanished.

The exact ranking representation also received a fair control. Ranking
with arity-coded terms has $2692/781/853$ in the scope-derived implementation,
whereas recursive formulas with the same term encoding have $2677/751/823$.
Thus rank flattening is a useful structural observation, but does not win
this comparison. The report selects the smaller verified implementation.

The selected packet protocol passed 620 standalone process checks, using normal Python and \texttt{python -O}. These include 86 accepted certificates, 86 exact packet round trips, and 52 rejection cases. A separate 2101 generated semantic checks compared opening with an independent finite numerical interpreter. There were zero failures. Additional tree-format and ranked-format suites check the intermediate translations; their separate results are included rather than added to the selected count.

Checks include hypotheses that are genuinely used, dependent paired
elimination, equation progress with a used induction hypothesis, witnesses
with nested iterator binders, malformed integers and lists, scope escapes,
forged Step premises, changed iterator bodies, missing clause premises,
and leftover packet input. The source control is frozen at the pinned
commit; its SHA-256 is
\begin{center}\ttfamily\scriptsize
17d116f7a6cbccaca169f76124aed5a9b20d45818747f96fba1c6f36668ad296.
\end{center}

\paragraph{What remains open.}
This work supplies a smaller complete checker in a transparent comparison
family and a four-scheme presentation with an explicit initial clause.
It does not establish a globally smallest axiomatization, does not improve
the repository's compact $516$-byte frontier, and does not assert novel
arithmetic strength. A future improvement should either delete genuinely
independent trusted structure or improve a matching complete implementation;
counting rule names or omitting typing and substitution costs is insufficient.

\section{A further logical quotient: rank from variable support}

There is a genuine logical quotient behind a tempting further simplification.
If the outer two variables of $D_{x,y}B$ are both absent from $B$, then
\[
 D_{x,y}B\ \leftrightarrow\ \exists x(B\to B)\ \leftrightarrow\ \top.
\]
For a ranked formula $(k,s,t)$, let $r$ be the largest free positive
de Bruijn index in the two endpoints, counted relative to iterator binders.
Take $r=0$ when the support is empty.
If $\lceil r/2\rceil<k$, its outer pair is vacuous and the formula can
be replaced by $0=0$. Otherwise $k=\lceil r/2\rceil$ and the public rank
can be reconstructed from the endpoints.

This does not remove ranks everywhere. A clause body has an additional
external universal slot at index $2k+1$. The same endpoint support can then
mean a rank-zero body with a clause variable or a public rank-one formula.
Internal clause ranks must still be retained unless the variable convention
is changed. The implemented support experiment grew to $3069$ normalized
bytes, with $854$ tokens and $957$ AST nodes. It is not selected.
Its source is supplied as an exploratory size comparison and is outside
the verified checker suite. The optional proof notes give the quotient
argument and the elimination case; the selected checker uses the exact
Eq/$D$ syntax and needs none of that extra machinery.

\section{Reproduction and artifact guide}

Run \verb|python3 reproduce.py| to execute the selected translations and
checks and regenerate measurements. Run
\verb|python3 reproduce.py --pdf| to compile the paper as well.
The selected checker runs directly on a certificate with only the Python
standard library. The accompanying README gives its exact path and sample
command, and records the compiler and verification environment.

The delivered \nolinkurl{report.tex} is self-contained: its checker listing
and inherited appendices are embedded. The source fragments and build
script are also supplied for editing. The inherited mathematical material
below is attributed to the pinned repository and is included to make the
arithmetic bridge inspectable offline. It is not presented as a new result
of this iteration.

\clearpage
\appendix
\section{Complete selected checker}
The following is the complete measured source. Imports, parsing, rejection,
scope checks, zero normalization, rule processing, and final acceptance are
all included.
\begin{lstlisting}[language=Python]
import json

def require(condition):
    if not condition:
        raise ValueError('Invalid certificate')

def term(a, scope, depth=0):
    if type(a) is int:
        return -scope <= a <= depth
    return type(a) is list and len(a) in (1, 3) and (len(a) == 1 or a[0] != 0) and all((term(t, scope, depth + (len(a) == 3 and i == 2)) for i, t in enumerate(a)))

def formula(a, scope, depth=0):
    return type(a) is list and len(a) == 2 and (formula(a[1], scope, depth + 2) if a[0] == 'd' else all((term(t, scope, depth) for t in a)))

def opening(a, value, depth=0):
    if type(a) is int:
        return value if a == depth + 1 else a - (a > depth + 1)
    if a[0] == 'd':
        return ['d', opening(a[1], value, depth + 2)]
    b = [opening(t, value, depth + (len(a) == 3 and i == 2)) for i, t in enumerate(a)]
    return b[1] if len(a) == 3 and b[0] == 0 else b

def check(goal, context=()):
    scope = len(context)
    require(formula(goal, scope))
    if goal[0] == goal[1]:
        return
    major, data = packets.pop()
    fresh = -scope - 1
    if type(major) is int:
        if major >= 0:
            require(major < len(context) and term(data, scope))
            antecedent, body = context[major]
            require(goal == opening(body, data))
            check(antecedent, context)
        else:
            require(major == -1 and goal[0] == 'd' and term(data, scope))
            body = opening(goal[1], data, 1)
            eta = fresh if goal[1][0] != 'd' and data == 0 else 0
            check(opening(body, [fresh]), context + (([0, 0], opening(body, eta)),))
    else:
        check(major, context)
        s, t = major
        if s == 'd':
            require(data == 0 and type(data) is int)
            body = opening(t, fresh, 1)
            check(goal, context + ((opening(body, 0), body),))
        else:
            check(data, context)
            match data:
                case [[[n], a, h], r]:
                    require(s == (a if n == 0 else [n, a, h]) and goal == [r, opening(h, t)])
                    return
            raise ValueError('Invalid step')
if __name__ == '__main__':
    packets = json.load(open(0))
    require(type(packets) is list and bool(packets))
    check(packets.pop())
    require(not packets)
    print('ACCEPT')
\end{lstlisting}

\clearpage
\section*{Inherited foundations: provenance}
\addcontentsline{toc}{section}{Inherited foundations: provenance}
The following five appendices reproduce the pinned
\nolinkurl{Research/certified_zero/reference/} foundations under the upstream
MIT-0 license. Their references to an ``old'' or ``new'' choice rule concern
that earlier basis. The identity-language equivalence in the present
Section~\ref{sec:equivalence} and the recovery in
Section~\ref{sec:recovery} connect that basis to the four-scheme calculus.
The appended proofs supply the logical interface, ordinary realization,
predecessor, equality reflection, and primitive-recursion compiler.

% New mathematical result: complementary introduction choices.
% Requires amsmath/amssymb/amsthm and theorem, lemma, proposition environments.
% The inherited arithmetic-to-iPRA bridge must be supplied separately.

\section{Complementary introductions with one local assumption}
\label{sec:choice-introduction}

The old paired introduction retains both the zero instance and an
induction instance when its matrix is an equation. Both together are
unnecessary. A certificate may choose either one, provided the induction
choice is restricted to a literal equation. This changes neither the
formula language nor the meaning of a formula. The omitted combination
is recovered by a short nested introduction--elimination construction.

\subsection{Complete native syntax and rule inventory}

The scalar numerical language is
\[
 t::=x\mid0\mid St\mid I(n,a;u.h).
\]
Only $h$ binds the state variable $u$; $n$ and $a$ do not.
The update may contain fixed numerical parameters. The computation
axioms in every retained context are
\[
 I(0,a;u.h)=a,\qquad
 I(Sn,a;u.h)=h[I(n,a;u.h)/u].
 \tag{I0, IS}\label{eq:choice-computations}
\]
The argument also applies unchanged to the full fixed-description PR
language with its projection, composition, and recursion equations.
No numerical variable ranges over function descriptions, and no uniform
evaluator is postulated.

Native formulas and reusable context clauses are
\[
 A::=(s=t)\mid\mathsf D_{x,y}A,\qquad
 \langle A;y.B(y)\rangle,\quad y\notin\operatorname{FV}(A).
\]
Bound variables are renamed before capture-free substitution.
The two variables $x,y$ in a paired formula are distinct.
An ordinary assumption $F$ is the constant clause
$\langle H;v.F\rangle$, where $v$ is absent from $F$ and $H$ is the
fixed computation theorem $I(0,0;u.0)=0$. Thus ordinary hypothesis
use is a clause Use preceded by this computation proof.

There are computation axioms and exactly the following inference
patterns:
\[
\frac{\Gamma\vdash A}{\Gamma\vdash B(t)}
\quad\mathsf{Use}
\quad(\langle A;y.B(y)\rangle\in\Gamma),
\]
\[
\frac{\Gamma\vdash\mathsf D_{x,y}B(x,y)\qquad
      \Gamma,\langle B(a,0);y.B(a,y)\rangle\vdash F}
     {\Gamma\vdash F}
\quad\mathsf{DE},
\qquad
\frac{\Gamma\vdash s=t\qquad\Gamma\vdash F(s)}
     {\Gamma\vdash F(t)}
\quad\mathsf{Tr},
\]
\[
\frac{\Gamma,B(t,\eta_\varepsilon(z))\vdash B(t,Sz)}
     {\Gamma\vdash\mathsf D_{x,y}B(x,y)}
\quad\mathsf{DI}_{\varepsilon},
\qquad
\eta_0(z)=0,\quad\eta_1(z)=z.
\tag{Choice DI}\label{eq:choice-di}
\]
Here $\varepsilon$ is a certificate choice in $\{0,1\}$, and
$\varepsilon=1$ is permitted only when $B$ has a literal equation
root. An equation beneath a binder is not a literal equation.
Call $\varepsilon=0$ the \emph{case choice} and $\varepsilon=1$
the \emph{progress choice}. Both retain exactly one ordinary
formula assumption. Compound matrices have only the case choice.

In DI the eigenparameter $z$ is fresh for $\Gamma$, the witness
$t$, and every fixed parameter of $B$. In DE the eigenparameter
$a$ is fresh for $\Gamma$, the conclusion $F$, and every fixed
parameter of $B$. In transport, $F(v)$ is any fixed native
formula template. Displayed local assumptions and clauses are
discharged. There are no primitive reflexivity, coverage,
separation, ordinary logical, or additional induction rules.

We call this calculus $\mathcal C$. The original comparison
calculus $\mathcal N$ has identical syntax, computation,
Use, DE, and transport, but its introduction is
\[
\frac{\Gamma,B(t,0),[B(t,z)]_{\mathrm{at}}\vdash B(t,Sz)}
     {\Gamma\vdash\mathsf D_{x,y}B(x,y)}.
\tag{Old DI}\label{eq:choice-old-di}
\]
The bracket supplies $B(t,z)$ exactly at a literal equation.

\subsection{Exact equivalence and the unused-witness construction}

\begin{lemma}[Structural admissibility]
\label{lem:choice-structural}
Weakening, exchange, contraction, renaming, ordinary-assumption cut,
and capture-free numerical substitution are admissible in $\mathcal C$.
\end{lemma}
\begin{proof}
Induct on derivations. Rename local eigenparameters before inserting
proofs or terms. Numerical substitution preserves a formula's root
constructor, so an allowed progress choice remains allowed.
The selected choice does not change under substitution or cut.
For ordinary-assumption cut, replace each use of its constant clause
by the supplied proof. Computation schemas are closed under
capture-free numerical substitution. These facts verify every case.
\end{proof}

\begin{theorem}[Identity-language equivalence]
\label{thm:choice-equivalence}
For every retained clause context $\Gamma$ and native formula $F$,
\[
 \Gamma\vdash_{\mathcal C}F
 \quad\Longleftrightarrow\quad
 \Gamma\vdash_{\mathcal N}F.
\]
The transformations are effective on finite derivations and leave
all numerical and formula syntax unchanged.
\end{theorem}
\begin{proof}
Each new introduction is an instance of the old one after weakening
by the omitted local assumption. For the case choice at an equation,
add the induction instance. For the progress choice, add the zero
instance. Compound case introduction is already the old rule.
Thus every new derivation translates to an old one.

For the reverse direction, transform an old derivation recursively.
Only its atomic DI case needs a construction. Suppose its matrix is
the equation $E(x,y)$ and its transformed premise is
\[
 \pi:\quad\Gamma,E(t,0),E(t,z)\vdash E(t,Sz).
\]
The outer new inference will be case DI with witness $t$. It is
therefore enough to prove $E(t,Sz)$ under
$\Gamma,E(t,0)$.

Rename the eigenvariable in $\pi$ to a fresh $v$. Choose a new
variable $u$ absent from $E(t,y)$ and all retained data. Under
$\Gamma,E(t,0)$, progress DI with witness zero introduces
\[
             \mathsf D_{u,y}E(t,y)
\]
from exactly the renamed premise $\pi[v/z]$.
The bound coordinate $u$ is deliberately unused.

Apply DE to this paired formula. Because $u$ is unused, the opened
clause is exactly
\[
               \langle E(t,0);y.E(t,y)\rangle;
\]
it contains no new existential witness. The retained ordinary
assumption $E(t,0)$ supplies its antecedent. Use at $Sz$ gives
$E(t,Sz)$, and DE closes. Finally the outer case DI introduces
the original $\mathsf D_{x,y}E(x,y)$ with witness $t$.
The original premise proof $\pi$ has been inserted once.

An old compound introduction already is the new case choice.
All other rule cases are unchanged. This completes the induction
on the original derivation.
\end{proof}

The unused inner coordinate is essential. Introducing
$\mathsf D_{x,y}E(x,y)$ internally and then eliminating it would
give an arbitrary fresh witness, not necessarily the original
$t$. The construction instead keeps $t$ as a fixed parameter
of $E(t,y)$ and binds only the unused coordinate $u$.

\begin{proposition}[Linear bound on inference nodes]
\label{prop:choice-node-bound}
Count each computation axiom, Use, DI, DE, and transport as one
inference node, expanding ordinary-assumption use as a computation
of $H$ followed by Use. If an old derivation has $N$ nodes and
$r$ atomic DI nodes, the displayed transformation has exactly
\[
                   N+5r\leq 6N
\]
nodes.
\end{proposition}
\begin{proof}
Each atomic old DI node is replaced by six nodes: outer case DI,
inner progress DI, DE, Use of the opened clause, Use of the
retained base assumption, and the computation of $H$ needed for
that last ordinary-assumption use. Its old premise subtree
occurs once. All other inference nodes are copied once.
\end{proof}

This counts inference-tree nodes. It is not a bound on serialized
JSON bytes, formula sizes, checker runtime, or the number of
consumed packet cells when hypothesis closure is implicit.
In the formula-context Q/D presentation, the corresponding
replacement has outer case DI, inner progress DI, DE, QE,
and two Hyp nodes, giving the same six-node count.

\subsection{Induction and numerical cases separate cleanly}

\begin{lemma}[Contextual equation induction]
\label{lem:choice-ei}
The progress choice, DE, and Use derive
\[
\frac{\Gamma\vdash E(0)\qquad
      \Gamma,E(z)\vdash E(Sz)}
     {\Gamma\vdash E(t)}
\tag{EI}\label{eq:choice-ei}
\]
for every equation $E$, with the usual freshness condition.
\end{lemma}
\begin{proof}
Choose an unused witness coordinate $u$. Progress DI introduces
$\mathsf D_{u,y}E(y)$ from the given step. DE opens the clause
$\langle E(0);y.E(y)\rangle$. Supply its antecedent using the given
base proof and use the clause at $t$. Close DE.
\end{proof}

\begin{lemma}[Cases, equality, and binder congruence]
\label{lem:choice-preinduction}
The case choice derives, for arbitrary native formulas,
\[
\frac{\Gamma\vdash F(0)\qquad\Gamma\vdash F(Sz)}
     {\Gamma\vdash F(t)}
\tag{NC}\label{eq:choice-nc}
\]
with the usual freshness condition. Equality reflexivity,
symmetry, transitivity, numerical congruence, and uniform
matrix interderivability under $\mathsf D$ are derivable
without using the progress choice.
\end{lemma}
\begin{proof}
For cases, use case DI with an unused witness coordinate and
matrix $F(y)$, ignoring its retained base in the successor
premise. DE opens $\langle F(0);y.F(y)\rangle$; the given zero
proof and Use at $t$ finish.

For reflexivity put $r=I(0,t;u.0)$. Use the computation equation
$r=t$ as both transport premises with motive $v=t$ to obtain
$t=t$. The usual explicit equality motives give symmetry,
transitivity, and numerical congruence.

For binder congruence, open the old paired formula at a fresh
witness $a$ and introduce the new one by case DI at $a$.
Convert its retained base to the old base, use the opened
old clause at the successor, and convert back. Ordinary cut
closes the construction, and the reverse conversion gives
the converse. No induction is used.
\end{proof}

The primitive progress choice is still restricted to equations.
Allowing it at arbitrary compound matrices would, by the
argument for \eqref{eq:choice-ei}, supply induction for
arbitrary formulas and change the intended theory.

\subsection{The entire ordinary logical interface}
\label{subsec:choice-logic}

Here is a predecessor-free reconstruction of the ordinary logical
operations. It makes explicit that the case choice supplies the
logical interface independently of progress induction.

Let $C(v,a,b):=I(v,a;u.b)$ with $u$ unused. Computation gives
$C(0,a,b)=a$ and $C(Sv,a,b)=b$. Define the gate on fixed finite
formula syntax by
\[
 G_v(s=t):=(C(v,0,s)=C(v,0,t)),\qquad
 G_v(\mathsf D_{x,y}B):=\mathsf D_{x,y}G_vB,
\]
renaming bound variables away from the control. Equality,
binder congruence, and external structural induction give
\[
 \vdash G_0A,\qquad G_{Sv}A\dashv\vdash A,\qquad
 A\dashv\vdash\mathsf D_{u,y}G_yA
 \quad(u,y\notin\operatorname{FV}(A)).
\tag{Gate and padding}\label{eq:choice-gate}
\]
For padding, case DI proves the forward direction. DE, the
zero-gate theorem, and Use at $1$ prove the reverse direction.

Define a selector between two arbitrary formulas by
\[
 M_v(s=t,u=w):=(C(v,s,u)=C(v,t,w))
\]
and, with aligned fresh paired variables, by
\[
\begin{aligned}
 M_v(\mathsf D_{x,y}A,\mathsf D_{x,y}B)
   &:=\mathsf D_{x,y}M_v(A,B),\\
 M_v(E,\mathsf D_{x,y}B)
   &:=\mathsf D_{x,y}M_v(G_yE,B),\\
 M_v(\mathsf D_{x,y}A,E)
   &:=\mathsf D_{x,y}M_v(A,G_yE),
\end{aligned}
\]
where $E$ is an equation. Every call decreases maximum input
binder depth. Padding and congruence prove
$M_0(A,B)\dashv\vdash A$ and
$M_{Sv}(A,B)\dashv\vdash B$.

Put $\delta(v)=C(v,1,0)$ and $p(z)=I(z,0;u.\delta(u))$.
Only the fixed computations $p(1)=1$ and $p(2)=0$ are needed.
The full logical macro inventory is
\[
\begin{aligned}
 A\Rightarrow_*B&:=\mathsf D_{u,z}M_z(A,B),\\
 A\wedge_*B&:=\mathsf D_{u,z}G_zM_{p(z)}(A,B),\\
 \exists_*x\,A(x)&:=\mathsf D_{x,y}G_yA(x),\\
 A\vee_*B&:=\exists_*x\,M_x(A,B),\\
 \forall_*x\,A(x)&:=
        A(0)\wedge_*\mathsf D_{u,z}G_zA(z),\\
 \bot_*&:=(0=1).
\end{aligned}
\tag{Logical macros}\label{eq:choice-macros}
\]
Newly bound variables are fresh; $u$ is unused.

Implication introduction obtains its assumption from the zero
selector and supplies the consequent at successors, then uses
case DI. Elimination opens its clause and uses it at $1$.
Existential introduction uses its witness and the successor
gate; elimination opens at a fresh witness, supplies the
zero-gate theorem, and uses the clause at $1$.

Disjunction introduction uses witnesses $0$ and $1$.
For elimination, discharge the two branch proofs as
$A\Rightarrow_*F$ and $B\Rightarrow_*F$. NC applied to
$M_z(A,B)\Rightarrow_*F$ gives this implication at every
control: the two endpoint cases follow from the branch
implications. Existential elimination provides $M_a(A,B)$
and the implication gives $F$.

Given $A$ and $B$, selector endpoints and NC give $M_v(A,B)$
at an arbitrary control. Substitute $v=p(Sz)$ and use the
successor gate for conjunction introduction. For projections,
open its zero-gated clause and specialize at $1$ and $2$;
the two fixed computations give $B$ and $A$, respectively.

These rules prove existential numerical coverage:
\[
     (x=0)\vee_*\exists_*v\,(x=Sv).
\tag{Coverage}\label{eq:choice-coverage}
\]
Use NC. Zero takes the left injection of reflexivity, and a
successor takes the right injection with the evident witness.
Coverage gives equality cases even when the retained context
contains $x$; its elimination adds $x=0$ or $x=Sv$ rather
than silently changing that retained context.

Universal introduction follows from a fresh-parameter proof
by substituting zero and successors and forming the two
conjuncts. For elimination, take its $A(0)$ conjunct and
open the paired conjunct. Its base is a zero-gate theorem.
Coverage of the target chooses the zero instance from
$A(0)$ or a successor instance from the opened clause
and gate removal.

Finally, under $0=1$, computation and transport give
$s=C(0,s,t)=C(1,s,t)=t$ for arbitrary numerical terms.
Every equation follows, and external induction on formula
syntax with case DI gives every native formula. This
supplies falsity elimination. No excluded middle for
quantified formulas has entered the construction.

\subsection{Whole-domain realizations and both formula composites}

For a standard first-order reading, let $T_{\rm eq}$ have
intuitionistic equality logic, the same numerical signature
and computation equations, numerical coverage, separation
$(0=1)\to\bot$, and contextual equation induction.
Its interpretation of the native syntax is
\[
 (s=t)^\circ=(s=t),\qquad
 (\mathsf D_{x,y}B)^\circ
   =\exists x\bigl(B^\circ(x,0)\to\forall y B^\circ(x,y)\bigr),
\]
\[
 \langle A;y.B\rangle^\circ
       =A^\circ\to\forall y B^\circ.
\]
The domain is the entire numerical domain and interpreted
equality is actual equality.

Case DI is sound in this realization by numerical coverage.
Progress DI is sound by contextual equation induction under
the conditional base. Use, DE, transport, and computation
have their displayed ordinary meanings. Conversely, the
native macros derive every ordinary logical rule,
\eqref{eq:choice-ei} derives equation induction, and
\eqref{eq:choice-coverage} derives coverage. Separation
translates to the identity implication
$(0=1)\Rightarrow_*(0=1)$.

For the essential paired formula composite,
\[
 \mathsf D_{x,y}B
 \dashv\vdash
 \exists_*x\bigl(B(x,0)\Rightarrow_*\forall_*y\,B(x,y)\bigr).
\tag{Paired composite}\label{eq:choice-paired-composite}
\]
Forwards, DE opens its clause at a fresh $a$; under $B(a,0)$
that clause supplies $B(a,y)$ at a fresh arbitrary $y$.
Introduce the universal, implication, and existential
macros and close DE. Backwards, eliminate the existential
at a fresh $a$. Introduce the paired formula with witness
$a$ by case DI: the retained base, supplied implication,
and universal elimination at $Sz$ give its successor premise.

The ordinary interpretations of the six logical macros
have their ordinary round trips. For conjunction, its true
base and the instances $1,2$ recover both operands;
the converse uses selector cases. The universal has all
successor instances in its paired conjunct and zero in
its first conjunct; coverage gives exactly the ordinary
universal. Implication, disjunction, and existential
quantification use their selector or gate endpoints.
Falsity uses the explicit ordinary separation bridge.
External induction on finite formulas, together with
\eqref{eq:choice-paired-composite}, gives both formula
composites. Clauses correspond in both directions to
$A\Rightarrow_*\forall_*yB$ by the derived logical rules.

The numerical composites are literal identities with
internally provable graph $x=y$. No standard-numeral
restriction, retract, quotient, or equality reflection
assumption is being inserted at this stage. The separately
supplied arithmetic bridge identifies $T_{\rm eq}$ with
full iPRA; the present change preserves that bridge
unchanged because Theorem~\ref{thm:choice-equivalence}
preserves every native sequent.

\subsection{What the theorem does and does not claim}

The new inference interface retains one ordinary assumption
at every paired introduction. It preserves the entire
consequence relation, not just equations or closed
theorems. Consequently the established full-iPRA
interpretations and equiconsistency transfer unchanged.

This result is different from retaining only progress
introduction and strengthening closure elimination.
Here the original elimination rule is retained, and
the case and progress introductions coexist. It is also
different from allowing arbitrary compound progress,
which would supply stronger induction.

A complete-checker comparison must still count the
certificate choice, its equation-root restriction,
substitution, context handling, computation, and the
public input/output interface. The mathematical theorem
alone is not a byte-count or minimality theorem.
The inference-node bound is distinct from source-size
measurements and from serialized certificate size.
The proof is a written universal proof; finite regression
tests do not constitute a Lean/Rocq formalization of it.


\clearpage
\section{Arithmetic strength and both ordinary interpretation composites}
\label{sec:ordinary-bridge}

We spell out the arithmetic bridge to avoid identifying standard numerical
agreement with bi-interpretability. First elaborate every fixed iterator
into an ordinary function symbol. At stage zero the symbols are $0,S$.
Given a fixed earlier-stage term $h(u,\vec p)$, introduce
$i_h(n,a,\vec p)$ with the two equations
\[
i_h(0,a,\vec p)=a,\qquad
i_h(Sn,a,\vec p)=h(i_h(n,a,\vec p),\vec p).
\]
Only finite descriptions occur; the domain contains no code for a
universal evaluator. Formal parameters are distinct when a description
is declared, while actual arguments can coincide.

Let $T_I$ be ordinary intuitionistic first-order logic with actual
equality and genuine falsity, these staged equations, contextual
induction for literal equations, and
\[
\forall x\,(x=0\vee\exists y\,x=Sy),
\qquad (0=S0)\to\bot.
\]
The second sentence is the falsity bridge. It must be stated: positive
numerical equations by themselves admit a singleton structure. Native
explosion from $0=S0$ and a translation of ordinary falsity to that
equation do not authorize silently dropping the bridge from $T_I$.

The native-to-ordinary translation is
\[
(s=t)^\circ=(s=t),\qquad
(\mathsf D_{x,y}B)^\circ=
\exists x\bigl(B^\circ(x,0)\to\forall y B^\circ(x,y)\bigr).
\]
A clause $\langle A;y.B\rangle$ becomes $A^\circ\to\forall yB^\circ$.
For the auxiliary closure used by the compact checkers,
$(\mathsf Q_yB)^\circ=B^\circ(0)\to\forall yB^\circ$.
The interpretation domain is all elements, defined by $x=x$;
equality is actual equality.

Case DI is sound because numerical coverage combines its retained base
with its successor proof. Progress DI is sound by literal-equation
induction under a temporarily assumed base. Use, DE and transport are
ordinary logical inferences under the translations. Conversely the
derived logical operations, equation induction, cases and explosion in
Section~\ref{sec:choice-introduction} recover the rules and axioms of
$T_I$. An auxiliary closure is interderivable with
$\mathsf D_{v,y}B(y)$ for unused $v$: in one direction use case DI with
the closure and its base; in the other eliminate the paired formula
and return its supplied closure. Thus formulas and context clauses have
the stated meanings in both presentations.

\subsection{The arithmetic dependency order}

The scalar update has no recursion-counter argument. We do not conceal
one in an encoding. The inherited internal proofs, reproduced in the
appendices, establish the required arithmetic in this order:
\begin{enumerate}
\item Using equation induction and numerical cases, construct a fixed
scalar predecessor term and prove $P0=0$ and $P(Sn)=n$.
Existential coverage then gives $x=0\vee x=S(Px)$: in the successor
case $x=Sy$, apply $P$ to obtain $Px=y$ and substitute.
\item Construct truncated subtraction and a numerical characteristic
of actual equality. Prove its reflection on the entire original domain.
\item Compile each fixed quantifier-free formula into its characteristic
equation and recover every contextual quantifier-free induction instance.
\item With that induction available, construct pairing and projections
by a bounded square-root search and compile every fixed general
primitive-recursive description into a scalar iterator term.
\end{enumerate}
No pairing or faithful equality detector is assumed to obtain
predecessor. No appeal to standard-model correctness replaces an
internal equality. Every induction fixes its passive parameters;
diagonal substitution occurs only into a previously proved identity.

For a target recursion
\[
f(0,\vec p)=g(\vec p),\qquad
f(Sn,\vec p)=h(n,f(n,\vec p),\vec p),
\]
let $\pi,L,H$ be the proved pairing and projections. The compiler uses
\[
V(n,\vec p)=I\bigl(n,\pi(0,g^*(\vec p));
u.\pi(SL(u),h^*(L(u),H(u),\vec p))\bigr),
\qquad f^*(n,\vec p)=H(V(n,\vec p)).
\]
Equation induction proves $L(V(n,\vec p))=n$; the two projection
identities then prove the target defining equations. The construction
is external recursion on each finite PR description, not an internal
quantification over arbitrary functions.

\begin{theorem}[Whole-domain bi-interpretability and equiconsistency]
The ordinary realization $T_I$ and full first-order iPRA are
bi-interpretable. In both theories $x=y$ defines the isomorphism from
the original domain to its double interpretation. The native Choice-DI
calculus is therefore a faithful proof presentation of a theory
bi-interpretable and equiconsistent with iPRA.
\end{theorem}
\begin{proof}
Interpret an iterator in iPRA by primitive recursion that ignores its
counter argument in the update. iPRA proves its equations, literal
equation induction, coverage and separation. In the reverse direction
use the scalar compiler. The predecessor and equality-reflection proofs
supply the successor axioms and every quantifier-free induction instance.
All ordinary logical rules are preserved.

For two iterators with the same seed and provably equal updates,
equation induction on the equality of their values proves uniqueness at
every counter, with all parameters fixed. The corresponding argument
applies to ordinary primitive recursion. Induct externally on a finite
description: its double translation satisfies its original defining
equations, so internal uniqueness identifies the two functions. This
also proves coherence when substitution changes the staged syntactic
description of an update.

The domain and equality are unchanged; hence the identity graph is
internally total, bijective and preserves every interpreted operation.
The displayed ordinary meanings of the logical macros and the paired
binder give both formula composites by external induction on finite
formula syntax. Clauses are similarly equivalent to their conditional
universal meanings. The explicit falsity bridge preserves contradiction
in both directions. Effective translations of finite proofs give
equiconsistency, without a claim that either theory proves its own
consistency.
\end{proof}


\clearpage
% Adapted with attribution from the pinned unparenthesized one-hole archive.
\section{Internal scalar predecessor bootstrap}
\label{sec:hybrid-bootstrap}

This scalar bootstrap is inherited from the supplied scalar development~\cite{erasure} and is reproduced to keep the scalar checker self-contained. It is not a new contribution of this paper. The proof uses a factor map fixing odd states and a saturated-iteration construction of predecessor.

The available arithmetic terms are $0$, $S$, and
\[
 I(n,a;u.h),\qquad
 I(0,a;u.h)=a,\qquad
 I(Sn,a;u.h)=h[I(n,a;u.h)/u].
\]
The update binds only $u$.  All other variables in it are fixed numerical
parameters.  We use equality reasoning, the already reconstructed
zero/successor cases, and induction on a single literal equation.
Every induction below keeps its other parameters fixed.  Predecessor,
successor injectivity, pairing, and induction on arbitrary
quantifier-free formulas are not assumptions of the argument.

\subsection{Elementary terms and two finite cycles}

Write $1=S0$ and use larger numerals as successor abbreviations.  Define
\[
 C(n,a,b)=I(n,a;u.b),\qquad
 Z(n)=C(n,1,0),\qquad
 a+b=I(b,a;u.Su).
\]
The update variable of $C$ is absent from its body.  Consequently
\[
 C(0,a,b)=a,\qquad C(Sn,a,b)=b.
\]
Numerical cases show that $Z(n)$ is $0$ or $1$.
The usual commutative-monoid equations for addition follow by literal
equation induction.  For example, $0+b=b$ and $Sa+b=S(a+b)$ are
proved by induction on $b$; these yield commutativity by a further
equation induction, and associativity follows by induction on the last
summand.  We use these equations to reorder finite sums.  The notation
$2a$ means $a+a$ and $3a$ means $a+a+a$; no multiplication operation
is needed for this bootstrap.

Put
\[
 N(n)=I(n,0;u.Z(u)),\qquad
 j(u)=Z(u)+2N(u),\qquad
 M(n)=I(n,0;u.j(u)).
\]
The range of $N$ is contained in $\{0,1\}$: its zero value is $0$,
and its successor value is a value of $Z$.  This uses numerical cases
on its counter, not induction on a disjunction.  Hence
\begin{equation}
 N(n)+Z(N(n))=1,\qquad
 Z(Z(N(n)))=N(n),\qquad N(n+2)=N(n).
 \label{eq:hybrid-parity}
\end{equation}
Here the last equation follows from two successor computations and
the preceding equation.  Literal-equation inductions give
\begin{equation}
 N(2n)=0,\qquad N(2n+1)=1,\qquad N(3n)=N(n).
 \label{eq:hybrid-parity-scaling}
\end{equation}
For the first equation the successor step adds two to the argument of
$N$; the second follows by one computation.  In the third, increasing
$n$ adds three to the left argument and therefore applies $Z$ once to
both sides of the induction equation.

Finite computations now give
\[
 j(0)=1,\qquad j(1)=2,\qquad j(2)=0.
\]
The proof below needs just these three values and
$M(Sn)=j(M(n))$; it does not assume a division algorithm or a theorem
characterizing $M$ as remainder modulo three.

\subsection{A ramp with a global scaling identity}

Define
\[
 G(n)=I(n,0;u.u+M(u)+1).
\]
The central invariant is
\begin{equation}
 M(G(n))=N(n).
 \label{eq:hybrid-ramp-invariant}
\end{equation}
Its zero case is immediate.  For its successor step, retain the
single induction equation and split the Boolean value $N(n)$.
If it is $0$, the next $G$ state is $G(n)+1$; its $M$ value is
$j(0)=1$, equal to $N(Sn)$.  If it is $1$, the next $G$ state
is $G(n)+2$; two computations of $M$ change its value from $1$
to $2$ and then to $0$, again equal to $N(Sn)$.  This proves
the invariant by literal-equation induction.

It follows that
\[
 G(Sn)=G(n)+N(n)+1,\qquad G(n+2)=G(n)+3.
\]
Equation induction therefore yields
\begin{equation}
 G(2n)=3n.
 \label{eq:hybrid-ramp-double}
\end{equation}
A separate literal-equation induction proves
\begin{equation}
 G(3n)=3G(n)+N(n).
 \label{eq:hybrid-ramp-triple}
\end{equation}
For completeness, its successor calculation is
\begin{align*}
 G(3Sn)
 &=G(3n)+N(3n)+4
   =3G(n)+2N(n)+4,\\
 3G(Sn)+N(Sn)
 &=3G(n)+3N(n)+3+Z(N(n)).
\end{align*}
The expressions agree in both Boolean cases for $N(n)$.
The base is $0=0$.  In particular,
\eqref{eq:hybrid-ramp-triple} was proved internally; it was not
inferred from an intended floor-function interpretation of $G$.

Set
\[
 H(u)=C(N(u),G(u),u).
\]
The crucial choice is to retain $u$ when its parity is $1$.
The preceding equations give
\begin{equation}
 H(2u)=3u,\qquad H(1)=1,\qquad H(3u)=3H(u).
 \label{eq:hybrid-factor-map}
\end{equation}
To check the last identity, use $N(3u)=N(u)$ and split the Boolean
$N(u)$.  At $0$, both sides are $3G(u)$ by
\eqref{eq:hybrid-ramp-triple}; at $1$, both are $3u$.
This argument does not require that an arbitrary number has already
been represented as $2v$ or $2v+1$.

\subsection{A comparator from scalar iteration}

For every fixed update $h$, equation induction on $n$, keeping $m,a$
and all parameters fixed, proves the semigroup identity
\[
 I(m+n,a;h)=I(n,I(m,a;h);h).
\]
The base is the zero computation and the step applies $h$ to the
induction equation. This identity will also be used in the later
pairing construction.


For a fixed update $h$, equation induction on $y$, with its initializer
$a$ and all parameters of $h$ fixed, proves the rotation identity
\begin{equation}
 I(Sy,a;h)=I(y,h(a);h).
 \label{eq:hybrid-rotation}
\end{equation}
Both sides have base $h(a)$, and their successor equations apply the
same $h$ to the previous equality.  Here and below $I(y,a;h)$ is
an abbreviation for one fixed finite scalar term, not quantification
over functions.

By \eqref{eq:hybrid-factor-map}, another equation induction gives
\begin{equation}
 I(y,3a;H)=3I(y,a;H).
 \label{eq:hybrid-iterator-scaling}
\end{equation}
Its base is $3a=3a$.  In the step the fixed-parameter induction
equation rewrites the new left side as
\[
 H(3I(y,a;H))=3H(I(y,a;H)).
\]
The initializer parameter stays $a$ throughout this induction.
Equation induction also gives $I(y,1;H)=1$, since $H(1)=1$.

Define
\[
 E(x)=I(x,1;u.2u),\qquad
 L(x,y)=N\bigl(I(y,E(x);H)\bigr).
\]
Then the following entire recursion specification is provable:
\begin{equation}
 L(0,y)=1,\qquad L(Sx,0)=0,\qquad L(Sx,Sy)=L(x,y).
 \label{eq:hybrid-comparator}
\end{equation}
The first equation uses $E(0)=1$ and the fixed-point identity for
$H$.  The second uses $E(Sx)=2E(x)$ and $N(2a)=0$.
For the third,
\begin{align*}
 L(Sx,Sy)
 &=N\bigl(I(Sy,2E(x);H)\bigr)\\
 &=N\bigl(I(y,H(2E(x));H)\bigr)\\
 &=N\bigl(I(y,3E(x);H)\bigr)\\
 &=N\bigl(3I(y,E(x);H)\bigr)\\
 &=N\bigl(I(y,E(x);H)\bigr)=L(x,y).
\end{align*}
The successive steps use rotation, $H(2a)=3a$, iterator scaling,
and $N(3a)=N(a)$.  Thus $L$ has the comparison equations at all
arguments internally.  Its standard-model meaning as the indicator
of $x\le y$ is motivation, not a premise of the following argument.

\subsection{Saturated iteration gives predecessor}

We now need only the three equations
\eqref{eq:hybrid-comparator}.  Define
\[
 B_m(u)=C(L(Su,m),u,Su),\qquad
 W(t,m)=I(t,0;u.B_m(u)).
\]
The subscript in $B_m$ records a fixed parameter; $B_m$ is an open
term abbreviation, not a function variable.

Numerical cases on $e$ prove the selector equation
\[
 C(e,Sa,Sb)=S C(e,a,b).
\]
Equivalently one may use equation induction on $e$, whose induction
hypothesis is unused.  Therefore
\begin{align}
 B_{Sm}(Su)
 &=C(L(SSu,Sm),Su,SSu)\notag\\
 &=C(L(Su,m),Su,SSu)=S B_m(u),
 \label{eq:hybrid-conjugacy}\\
 B_{Sm}(0)&=C(L(S0,Sm),0,S0)=C(L(0,m),0,S0)=S0.
 \label{eq:hybrid-conjugacy-base}
\end{align}

Literal-equation induction on $t$, with $m$ fixed, proves
\begin{equation}
 W(St,Sm)=S W(t,m).
 \label{eq:hybrid-saturated-shift}
\end{equation}
The base is \eqref{eq:hybrid-conjugacy-base}.  The successor step is
\begin{align*}
 W(SSt,Sm)
 &=B_{Sm}(W(St,Sm))\\
 &=B_{Sm}(S W(t,m))\\
 &=S B_m(W(t,m))=S W(St,m).
\end{align*}
Only the second line uses the current induction hypothesis, at the
same fixed $m$.

Two subsequent equation inductions give
\begin{equation}
 L(Sn,n)=0,\qquad W(n,n)=n.
 \label{eq:hybrid-diagonals}
\end{equation}
The first uses the comparator recurrence.  The second uses
\eqref{eq:hybrid-saturated-shift} in its successor step.  The
substitution $t=m=n$ is made into that already proved universally
parameterized equation; no parameter of an active induction hypothesis
is changed.  Consequently
\begin{equation}
 W(Sn,n)=B_n(n)=C(L(Sn,n),n,Sn)=n.
 \label{eq:hybrid-extra-step}
\end{equation}

Finally define the finite scalar term
\begin{equation}
 P(m)=I\bigl(m,0;u.C(L(SSu,m),u,Su)\bigr).
 \label{eq:hybrid-predecessor-term}
\end{equation}
Its zero equation is immediate.  To justify its successor equation
without assuming a separate extensionality rule for iterator bodies,
put
\[
 V(t,m)=I\bigl(t,0;u.C(L(SSu,Sm),u,Su)\bigr).
\]
The comparator recurrence identifies this update pointwise with
$B_m(u)$.  Equation induction on $t$, with $m$ fixed, gives
$V(t,m)=W(t,m)$: their base values are $0$, and the successor step
uses the pointwise update equation and then the induction equation.
Now \eqref{eq:hybrid-extra-step} yields
\[
 \boxed{\qquad P(0)=0,\qquad P(Sm)=V(Sm,m)=W(Sm,m)=m.\qquad}
\]

\subsection{Dependency boundary}

All terms in this section are built from $0,S,I$ by finite substitution.
All inductions have a single equation as their invariant.  Numerical
cases and equality transport are the already reconstructed logical
resources.  Successor injectivity may now be obtained by applying $P$
to an assumed equation $Sx=Sy$.  Equality tests, reflection of
quantifier-free formulas, pairing, and the recovery of general
primitive recursion can then be developed in the established order.
They were not used in order to manufacture predecessor.

This argument is a written internal derivation in the stated theory.
The finite implementation checks do not constitute a proof-assistant
formalization of this entire metatheorem.


\clearpage
% Standalone section.  Requires only amsmath and amssymb in the parent document.
% No theorem environments, macros, or bibliography keys are assumed.
% All arithmetic claims in this section are proved here.

\section{From literal-equation induction to full intuitionistic PRA}
\label{sec:eic-completion}\label{app:arithmetic}

This section supplies an unconditional completion theorem. It is adapted
with attribution from the current forward-left draft~\cite{forward}. In the
scalar application used here, every function occurring in the proof is
already a scalar term: predecessor is supplied by
Appendix~\ref{sec:hybrid-bootstrap}, the selector is
$C(n,a,b)=I(n,a;u.b)$, addition is $x+y=I(y,x;u.Su)$, and
truncated subtraction is $d(x,y)=I(y,x;u.Pu)$. All other functions below
are finite compositions of these. Thus this proof does not presuppose
that arbitrary primitive recursion has already been compiled. The QFI
conclusion first applies to the scalar language; the general PR compiler
comes afterward.  Literal-equation
induction, together with the primitive-recursive defining equations,
predecessor coverage, and separation of $0$ and $1$, proves induction for every
quantifier-free formula.  In particular, the proof does not assume a faithful
numerical equality test, stability or decidability of equality, successor
uniqueness, or any commutation law for an iterator.  We construct the equality
test and prove its reflection property for arbitrary elements of the
numerical domain.

\subsection{The exact hypotheses}

Work in ordinary intuitionistic first-order logic with equality and genuine
falsity $\bot$.  The language has one numerical sort and a symbol for each
fixed, finite primitive-recursive description.  Distinct descriptions are
not identified merely because they compute the same standard function.
Let $\mathsf{PRE}$ contain the defining equations for zero, successor,
projections, composition, and primitive recursion.  In particular, choose
descriptions $P$ and $C$ with equations
\begin{align}
 P(0)&=0, & P(Sx)&=x, \label{eq:eic-predecessor}\\
 C(0,a,b)&=a, & C(Sx,a,b)&=b. \label{eq:eic-selector}
\end{align}
Write $1=S0$.  In addition to $\mathsf{PRE}$, assume
\begin{align}
 \mathsf{Cov}:&\qquad
       \forall x\,\bigl(x=0\lor x=S(Px)\bigr),
       \label{eq:eic-coverage}\\
 \mathsf{Sep}:&\qquad 0=1\longrightarrow\bot,
       \label{eq:eic-separation}
\end{align}
and induction for each \emph{literal equation} $E(n,\vec p)$:
\begin{equation}
 \forall\vec p\,
 \left[
   \left(E(0,\vec p)\land
       \forall n\bigl(E(n,\vec p)\longrightarrow E(Sn,\vec p)\bigr)
   \right)
   \longrightarrow\forall n\,E(n,\vec p)
 \right].
 \label{eq:eic-induction}
\end{equation}
Here the schematic expression $E$ is a single equality of numerical
terms; the displayed universal quantifiers range only over numbers.
Every description and every instance of the schema is finite external
syntax.

The axiom formulation \eqref{eq:eic-induction} and the usual contextual
induction rule are interchangeable.  For example, if
\[
 \Gamma\vdash E(0,\vec p),
 \qquad
 \Gamma,E(z,\vec p)\vdash E(Sz,\vec p),
\]
where $z$ is absent from $\Gamma$ and the passive parameters, discharge the
second premise's equation and universally introduce $z$.  The axiom then
gives $\Gamma\vdash\forall n E(n,\vec p)$, which can be instantiated at any
target term.  Conversely, assume the antecedent of the axiom and apply the
contextual rule with fresh variables.  This uses the ordinary intuitionistic
logical rules and permits arbitrary, including quantified, formulas in
$\Gamma$.

\subsection{Constructor separation and two elementary tests}

Define the sign description and the countdown description by
\begin{equation}
 \operatorname{sg}(u)=C(u,0,1),
 \qquad
 K(u)=C(u,1,Pu).
 \label{eq:eic-sign-countdown}
\end{equation}
Their defining equations imply
\begin{align}
 \operatorname{sg}(0)&=0,
 &\operatorname{sg}(Sx)&=1,\label{eq:eic-sign-computation}\\
 K(0)&=1,
 &K(Sx)&=x.\label{eq:eic-countdown-computation}
\end{align}

\paragraph{Constructor laws.}
Successor is injective: applying $P$ to $Sx=Sy$ gives $x=y$ by
\eqref{eq:eic-predecessor}.  Successor never equals zero.  Indeed,
$Sx=0$ gives $1=0$ by applying $\operatorname{sg}$ and using
\eqref{eq:eic-sign-computation}; symmetry and
\eqref{eq:eic-separation} give $\bot$.  Consequently $1\ne0$, where
$A\ne B$ abbreviates $A=B\to\bot$.  None of these arguments uses
induction.

\paragraph{The sign test.}
For every $u$,
\begin{align}
 &\operatorname{sg}(u)=0\lor\operatorname{sg}(u)=1,
       \label{eq:eic-sign-range}\\
 &\operatorname{sg}(u)=0\quad\longleftrightarrow\quad u=0,
       \label{eq:eic-sign-zero}\\
 &\operatorname{sg}(u)=1\quad\longleftrightarrow\quad u\ne0.
       \label{eq:eic-sign-one}
\end{align}
To prove the range assertion, split \eqref{eq:eic-coverage}.  The branch
$u=0$ gives sign $0$, and the branch $u=S(Pu)$ gives sign $1$.
The implication $u=0\to\operatorname{sg}(u)=0$ is computation.
For its converse, the successor branch of coverage would make
$\operatorname{sg}(u)=0$ imply $1=0$, a contradiction; in that branch
explosion gives $u=0$ as well.  Similarly, if
$\operatorname{sg}(u)=1$ and $u=0$, computation gives $0=1$.
Conversely, $u\ne0$ excludes the zero branch of coverage, and the
successor branch computes the sign as $1$.

Coverage and \eqref{eq:eic-selector} also give, for all $u,v$,
\begin{equation}
 C(u,v,v)=v.
 \label{eq:eic-selector-idempotence}
\end{equation}
Both coverage branches compute to $v$, so this identity again needs no
induction.

\paragraph{The countdown test and the notation $W$.}
For every $u$,
\begin{equation}
 K(u)=0\quad\longleftrightarrow\quad u=1.
 \label{eq:eic-countdown-zero}
\end{equation}
If $u=1=S0$, \eqref{eq:eic-countdown-computation} gives $K(u)=0$.
Conversely, assume $K(u)=0$ and split coverage.  If $u=0$, then
$K(u)=1$, contradicting separation, and explosion gives $u=1$.
If $u=S(Pu)$, computation gives $K(u)=Pu$; hence $Pu=0$ and
$u=S0=1$.

If a presentation uses a description named $W$ with defining equations
\[
 W(0)=1,\qquad W(Sx)=x,
\]
then it proves $K(u)=W(u)$ for every $u$ already from coverage.
In the zero branch both sides are $1$; in the branch $u=S(Pu)$ both
sides are $Pu$.  Thus $K$ and $W$ are provably equal on the whole domain,
even if they are distinct descriptions.  This equality does not appeal to
standard-model agreement or to uniqueness of an arbitrary recurrence.

\subsection{Four elementary uses of equation induction}

Choose descriptions for addition and truncated subtraction, with recursion
on the second argument:
\begin{align}
 x+0&=x, &x+Sy&=S(x+y),\label{eq:eic-addition}\\
 d(x,0)&=x, &d(x,Sy)&=P(d(x,y)).\label{eq:eic-subtraction}
\end{align}
Using these equations and \eqref{eq:eic-induction}, we now derive
\begin{align}
 0+y&=y,\label{eq:eic-left-zero}\\
 (Sx)+y&=S(x+y),\label{eq:eic-left-successor}\\
 d(Sx,Sy)&=d(x,y),\label{eq:eic-cancel-successors}\\
 d(Sa,a)&=1.\label{eq:eic-diagonal}
\end{align}
Each induction below has a single equation as its induction formula.

For \eqref{eq:eic-left-zero}, induct on $y$.  The base is $0+0=0$.
The step is $0+Sy=S(0+y)=Sy$.

For \eqref{eq:eic-left-successor}, hold $x$ fixed and induct on $y$.
At zero both sides are $Sx$.  At a successor the induction hypothesis
gives
\[
 (Sx)+Sy=S\bigl((Sx)+y\bigr)
       =S\bigl(S(x+y)\bigr)=S(x+Sy).
\]
Combining \eqref{eq:eic-left-successor} with the defining equation for
addition gives the useful rearrangement
\begin{equation}
 (Sy)+r=y+Sr.
 \label{eq:eic-shift-successor}
\end{equation}
Indeed, both sides are $S(y+r)$.

For \eqref{eq:eic-cancel-successors}, hold $x$ fixed and induct on $y$.
The base computes as
\[
 d(Sx,S0)=P(d(Sx,0))=P(Sx)=x=d(x,0).
\]
Applying $P$ to the induction equation gives the step:
\[
 d(Sx,S(Sy))=P(d(Sx,Sy))=P(d(x,y))=d(x,Sy).
\]

Finally, induct on $a$ for \eqref{eq:eic-diagonal}.  The base is
$d(S0,0)=S0=1$.  At a successor,
\[
 d(S(Sa),Sa)=d(Sa,a)=1,
\]
using \eqref{eq:eic-cancel-successors} and the induction hypothesis.

\subsection{An equation invariant that reflects actual equality}

The remaining arithmetic ingredient is the following equation, valid for
all $x,y$:
\begin{equation}
 C\bigl(\operatorname{sg}(d(x,y)),\ x,\ y+d(x,y)\bigr)=x.
 \label{eq:eic-residual-invariant}
\end{equation}
Its meaning is that a nonzero subtraction residual reconstructs the
original argument by addition.  The selector expresses this assertion as
one literal equation, so equation induction applies directly.

Hold $x$ fixed and induct on $y$.  At $y=0$, subtraction gives residual
$x$, and \eqref{eq:eic-left-zero} gives $0+x=x$.  The desired equation
is consequently $C(\operatorname{sg}(x),x,x)=x$, which is
\eqref{eq:eic-selector-idempotence}.

For the step, assume \eqref{eq:eic-residual-invariant} at $y$ and
abbreviate
\[
 r=d(x,y),\qquad r'=Pr=d(x,Sy).
\]
These are local term abbreviations, not additional function symbols.
Split coverage for $r'$.

If $r'=0$, then $\operatorname{sg}(r')=0$, and the successor instance
of \eqref{eq:eic-residual-invariant} reduces immediately to $x=x$.

In the other branch, $r'=S(Pr')$, so constructor separation gives
$r'\ne0$.  Split coverage for $r$.  Its zero branch would imply
$r'=P0=0$, a contradiction, from which the desired step equation follows.
In its successor branch,
\[
 r=S(Pr)=Sr',\qquad \operatorname{sg}(r)=1.
\]
The induction equation therefore computes to $y+r=x$.  Meanwhile,
\eqref{eq:eic-shift-successor} gives
\[
 (Sy)+r'=y+Sr'=y+r=x.
\]
Since $r'$ is a successor, $\operatorname{sg}(r')=1$, so this is
exactly the branch selected in the successor instance of
\eqref{eq:eic-residual-invariant}.  This completes the equation induction.
All case distinctions came from coverage; no induction on an implication
or a disjunction was used.

\paragraph{Reflection at residual $1$.}
We have, for arbitrary $a,b$,
\begin{equation}
 d(Sa,b)=1\quad\longleftrightarrow\quad a=b.
 \label{eq:eic-residual-reflection}
\end{equation}
For the forward direction, substitute $x=Sa$, $y=b$ into
\eqref{eq:eic-residual-invariant}.  If $d(Sa,b)=1$, the sign is $1$
and the selector returns its second branch, giving
\[
 b+1=Sa.
\]
The defining equation for addition gives $b+1=S(b+0)=Sb$.  Thus
$Sb=Sa$, and applying $P$ gives $b=a$, hence $a=b$.
Conversely, $a=b$ and \eqref{eq:eic-diagonal} give
$d(Sa,b)=d(Sa,a)=1$.

Define the equality characteristic by the fixed description
\begin{equation}
 e(a,b)=\operatorname{sg}\bigl(K(d(Sa,b))\bigr).
 \label{eq:eic-equality-characteristic}
\end{equation}
The sign, countdown, and residual reflection lemmas give
\begin{align}
 e(a,b)=0&\quad\longleftrightarrow\quad a=b,
       \label{eq:eic-equality-reflection}\\
 e(a,b)=0&\ \lor\ e(a,b)=1.
       \label{eq:eic-equality-range}
\end{align}
In detail, $e(a,b)=0$ is equivalent to $K(d(Sa,b))=0$ by
\eqref{eq:eic-sign-zero}, which is equivalent to $d(Sa,b)=1$ by
\eqref{eq:eic-countdown-zero}, and then to $a=b$ by
\eqref{eq:eic-residual-reflection}.  The range assertion is
\eqref{eq:eic-sign-range} applied to $K(d(Sa,b))$.

The reflection property concerns the logic's actual equality throughout
the domain.  It was not assumed as part of a numerical comparison symbol,
and its proof did not restrict $a$ or $b$ to standard numerals.
It also yields decidable equality.  In the $e(a,b)=0$ branch,
\eqref{eq:eic-equality-reflection} gives $a=b$.  In the $e(a,b)=1$
branch, an assumption $a=b$ would give $e(a,b)=0$ and therefore $0=1$,
contradicting separation.  Hence
\begin{equation}
 \forall a\,\forall b\,\bigl(a=b\lor(a=b\to\bot)\bigr)
 \label{eq:eic-equality-decidable}
\end{equation}
is a theorem, not an extra hypothesis.

\subsection{Compilation of all quantifier-free formulas}

For each fixed quantifier-free formula $A(\vec p)$, construct a term
$b_A(\vec p)$ by external recursion on its finite syntax.  Take equality,
$\bot$, $\land$, $\lor$, and $\to$ as the primitive formula constructors;
negation is implication into $\bot$.  Set
\begin{align}
 b_{s=t}&=e(s,t), & b_{\bot}&=1,\label{eq:eic-boolean-atoms}\\
 b_{A\land D}&=C(b_A,b_D,1),\label{eq:eic-boolean-and}\\
 b_{A\lor D}&=C(b_A,0,b_D),\label{eq:eic-boolean-or}\\
 b_{A\to D}&=C(b_A,b_D,0).\label{eq:eic-boolean-implies}
\end{align}
Every term here is a finite composition of fixed primitive-recursive
descriptions; if desired, projections and composition give it a single
description symbol.  No internal quantification over formulas or programs
is involved.

An external induction on the syntax of $A$ proves, within the arithmetic
theory,
\begin{equation}
 \bigl(b_A=0\lor b_A=1\bigr)
 \quad\text{and}\quad
 \bigl(b_A=0\longleftrightarrow A\bigr).
 \label{eq:eic-boolean-correctness}
\end{equation}
Here are the full connective checks.  In every case $b_A=1$ implies
$\neg A$: if $A$ held, the already established implication
$A\to b_A=0$ would contradict separation.

\paragraph{Atomic formulas and falsity.}
For an equation, both claims are
\eqref{eq:eic-equality-reflection}--\eqref{eq:eic-equality-range}.
For falsity the range assertion is $1=0\lor1=1$, with its right disjunct
proved by reflexivity.  Also $1=0\to\bot$ follows from separation and
symmetry, while $\bot\to1=0$ follows by explosion.  Thus $1=0$ is
equivalent to $\bot$, as required.

\paragraph{Conjunction.}
Split $b_A=0\lor b_A=1$.  In the zero branch, $A$ holds and
$b_{A\land D}=b_D$.  Its range follows from the range of $b_D$,
and its being zero is equivalent to $D$, hence to $A\land D$ under
the known $A$.  In the one branch, $\neg A$ holds and
$b_{A\land D}=1$.  Its range is immediate.  If it were zero,
separation would give $\bot$, hence $A\land D$; conversely an
assumption $A\land D$ contradicts $\neg A$, and explosion gives
$b_{A\land D}=0$.

\paragraph{Disjunction.}
In the branch $b_A=0$, the formula $A$ holds and
$b_{A\lor D}=0$.  The range is immediate, and both $b_{A\lor D}=0$
and $A\lor D$ hold, proving their equivalence.  In the branch
$b_A=1$, we have $\neg A$ and $b_{A\lor D}=b_D$.  Its range follows
from the range of $b_D$.  If its value is zero, then $D$ and hence
$A\lor D$.  Conversely, from $A\lor D$, the $A$ alternative
contradicts $\neg A$ and yields $b_D=0$ by explosion, while the
$D$ alternative yields $b_D=0$ directly.  This proves the required
equivalence in both alternatives.

\paragraph{Implication.}
In the branch $b_A=0$, the formula $A$ holds and
$b_{A\to D}=b_D$.  Its range follows from the range of $b_D$.
If its value is zero, then $D$, which gives $A\to D$.  Conversely,
$A\to D$ and the known $A$ give $D$, hence $b_D=0$.
In the branch $b_A=1$, we have $\neg A$ and
$b_{A\to D}=0$.  Its range is immediate.  To prove $A\to D$,
assume $A$; this contradicts $\neg A$, and explosion gives $D$.
Thus both $b_{A\to D}=0$ and $A\to D$ hold in this branch, and
their equivalence follows.  This argument uses the proved two-value
range of the finite characteristic.  It does not use excluded middle
for arbitrary formulas or change the surrounding intuitionistic logic.

This completes the external syntactic induction proving
\eqref{eq:eic-boolean-correctness} for every quantifier-free formula.

\subsection{Recovery of induction and the full theory}

Let $A(n,\vec p)$ be any fixed quantifier-free formula.  Suppose an
arbitrary retained context $\Gamma$ proves its base and its induction
step:
\[
 \Gamma\vdash A(0,\vec p),
 \qquad
 \Gamma,A(z,\vec p)\vdash A(Sz,\vec p),
\]
with the usual freshness condition on $z$.  By
\eqref{eq:eic-boolean-correctness}, the base gives
$\Gamma\vdash b_A(0,\vec p)=0$.  Under the additional assumption
$b_A(z,\vec p)=0$, the same equivalence gives $A(z,\vec p)$,
the given step gives $A(Sz,\vec p)$, and the equivalence gives
$b_A(Sz,\vec p)=0$.  Apply literal-equation induction to
$b_A(n,\vec p)=0$ and then apply its reflection implication.
The result is $\Gamma\vdash A(t,\vec p)$ for every target term $t$.
Thus the theory proves the full contextual quantifier-free induction
rule, including in arbitrary quantified contexts.

\paragraph{Completion theorem.}
In the common language of fixed primitive-recursive descriptions, the
theory
\[
 \mathsf{PRE}+\mathsf{Cov}+\mathsf{Sep}+\mathsf{EI}
\]
has exactly the same sequents as ordinary first-order intuitionistic
primitive recursive arithmetic with all quantifier-free induction
instances.  The forward simulation follows from the recovered
quantifier-free induction, successor injectivity, and successor separation;
all primitive-recursive defining equations and all logical rules are
shared.  For the reverse simulation, literal-equation induction is an
instance of quantifier-free induction, and successor separation gives
$\mathsf{Sep}$.  Quantifier-free induction proves $\mathsf{Cov}$:
at zero choose its left disjunct, and at a successor use
$Sx=S(P(Sx))$, which follows from $P(Sx)=x$, for its right disjunct.

Both interpretations therefore use the full domain, defined by $x=x$,
and preserve every function symbol, equality, logical connective, and
quantifier.  Both composites have the identity isomorphism, with graph
$x=y$.  The theories are consequently identity-domain bi-interpretable
and equiconsistent.  This completion theorem is unconditional once the
literal-equation induction rule has been obtained; it has no dependency
on a classical successor-uniqueness reduction.


\clearpage
\section{Pairing, primitive recursion, and internal numerical composites}
\label{sec:iterator-coding}
This appendix is adapted with attribution from the progress-closure
arithmetic development~\cite{progress}. Its prerequisites have now been
proved: scalar predecessor, arithmetic reflection, and full quantifier-free
induction. Pairing is a derived proof tool; it is not part of any checker.

\subsection{Pairing from a bounded square-root search}

Once QFI is available, a familiar square-based injection gives a
particularly direct compiler. Write $a\le b$ for $d(a,b)=0$.
The usual linear-order and semiring identities follow by fixed
quantifier-free inductions from the successor, addition, and monus
equations. Multiplication itself is the scalar term
\[
                        a\cdot b=I(b,0;u.u+a).
\]
Only primitive-recursive arithmetic is involved in the following
bounded root search; there is no minimization over an unbounded domain.

For a fixed parameter $n$, put
\[
 U_n(v)=C(d((Sv)^2,n),Sv,v),\qquad
 R_n(i)=I(i,0;v.U_n(v)),\qquad q(n)=R_n(n).
\]
The update increments $v$ exactly when $(v+1)^2\le n$, and
otherwise retains it. QFI in $i$, with $n$ fixed, proves
\[
 R_n(i)\le i,\qquad R_n(i)^2\le n,\qquad
 R_n(i)=i\ \vee\ n<(S(R_n(i)))^2.              \tag{Root invariant}
\]
The base has $R_n(0)=0$. If the update increments, its test
contradicts the strict-upper-bound alternative in the third invariant.
Thus the old state equals $i$, and the new state equals $Si$;
its square bound is the increment test. If the update does not increment, the
strict upper square bound holds and the state remains fixed.
An already established strict upper square bound forces every
subsequent update to remain fixed. These are exactly the cases
needed for the three displayed invariants.

At $i=n$, either the strict upper square bound already holds or
$R_n(n)=n$, in which case $n<(n+1)^2$ supplies it. Thus
\[
                        q(n)^2\le n<(q(n)+1)^2.           \tag{Root}
\]
Monotonicity of squares shows uniqueness of a number satisfying
Root. All formulas just used are quantifier free in the scalar
language, with fixed $n$ during the induction.

Now define
\[
 \pi(a,b)=(a+b)^2+a,\qquad
 L(n)=d(n,q(n)^2),\qquad H(n)=d(q(n),L(n)).
\]
For $k=a+b$, one has $k^2\le k^2+a<(k+1)^2$ since $a\le k$.
Root uniqueness gives $q(\pi(a,b))=a+b$, and therefore
\[
                    L(\pi(a,b))=a,\qquad H(\pi(a,b))=b.
\]
Surjectivity of $\pi$ is unnecessary. The original numerical domain
is never replaced by the range of this coding; pairing is used only
to store the finite state of a fixed primitive recursion. This
standard square pairing is a convenient reconstruction here, not
a claimed new arithmetic basis.

\subsection{Compiling each primitive recursion}

Suppose the fixed target description is
\[
 f(0,\vec p)=g(\vec p),\qquad
 f(Sn,\vec p)=h(n,f(n,\vec p),\vec p).
\]
Assume by external structural induction that the earlier $g,h$
already have scalar translations $g^*,h^*$. Define
\[
 W(u,\vec p)=\pi(SL(u),h^*(L(u),H(u),\vec p)),
\]
\[
 V(n,\vec p)=I(n,\pi(0,g^*(\vec p));u.W(u,\vec p)),
 \qquad f^*(n,\vec p)=H(V(n,\vec p)).
\]
The update has only its state variable bound, with every $p_i$
fixed. Equation induction proves $L(V(n,\vec p))=n$.
The two projection equations and scalar computation then prove
both defining equations for $f^*$ internally. Initial functions
and composition translate directly. This completes the translation
of every fixed finite PR description without adding a universal
evaluator.

\subsection{Why both numerical round trips are internal}

Two fixed scalar iterators with the same seed and provably equal
updates are internally equal at all counters: use literal-equation
induction comparing their values. Their base agrees, and their
successor equations plus congruence prove the step. The same
uniqueness argument works for ordinary primitive recursion.

For an ordinary PR description, induct externally on its finite
construction. The double translation satisfies its original
defining equations, so internal equation induction identifies it
with the original operation. For a native scalar description,
earlier translated updates agree by the external induction
hypothesis; scalar uniqueness identifies the double-translated
iterator. These proofs hold at arbitrary internal counters.

Thus the domain relation $x=y$ is provably a total bijection
preserving every fixed numerical operation in both composites.
The formula constructions and round trips in
Section~\ref{sec:choice-introduction} give the corresponding formula
composites. Both interpretations preserve falsity, taking native
$0=1$ to an ordinary contradiction and ordinary $\bot$ to native
$0=1$. Finite proof translations therefore establish equiconsistency.



\clearpage
\begin{thebibliography}{10}
\raggedright
\bibitem{baseline}
ProveIt contributors. \emph{Concise presentations of intuitionistic primitive
recursive arithmetic}, README at commit
\texttt{443cff49563457a48ee0f663012d56c4bdb58fa0}.
\url{https://github.com/VladimirReshetnikov/ProveIt/tree/443cff49563457a48ee0f663012d56c4bdb58fa0/Logic/PrimitiveRecursiveArithmetic/ConciseIPRA}.
All repository references below use this exact commit.

\bibitem{uniformstep}
ProveIt contributors. \emph{Uniform Step},
\nolinkurl{Research/uniform_step/}: report, claims, complete checker, and fixtures.
The frozen matched source is \nolinkurl{code/checker.py}.

\bibitem{compact}
ProveIt contributors. \emph{Canonical recursors} and \emph{Checked refinement},
\nolinkurl{Research/canonical_recursors/} and
\nolinkurl{Research/checked_refinement/}. Compact source-size comparisons and
their documented input contracts.

\bibitem{nativezerostep}
ProveIt contributors. \nolinkurl{Lean/ConciseIPRA/NativeZeroStep.lean}.
Declarations \texttt{NSPrf}, \texttt{nsprf\_iff\_nprf},
\texttt{nsprf\_iff\_pra}, and \texttt{NSPrf.normal}.

\bibitem{biinterp}
ProveIt contributors. \nolinkurl{Lean/ConciseIPRA/BiInterpretation.lean}.
Declarations \texttt{biInterp}, \texttt{teq\_roundTrip},
\texttt{pra\_roundTrip}, and \texttt{teq\_consistent\_iff}.

\bibitem{nativebi}
ProveIt contributors. \nolinkurl{Lean/ConciseIPRA/NativeBiInterpretation.lean}.
Declarations \texttt{native\_iff\_pra}, \texttt{pra\_iff\_native},
\texttt{native\_roundTrip}, \texttt{pra\_native\_roundTrip}, and
\texttt{native\_inconsistent\_iff}.

\bibitem{erasure}
ProveIt contributors. Inherited scalar predecessor bootstrap,
\nolinkurl{Research/certified_zero/reference/scalar_predecessor.tex}.
Reproduced in the inherited appendix.

\bibitem{forward}
ProveIt contributors. Inherited arithmetic reflection,
\nolinkurl{Research/certified_zero/reference/equality_reflection.tex}.
Reproduced in the inherited appendix.

\bibitem{progress}
ProveIt contributors. Inherited primitive-recursion compiler,
\nolinkurl{Research/certified_zero/reference/primitive_recursion.tex}.
Reproduced in the inherited appendix, together with
\nolinkurl{inherited_choice_basis.tex} and
\nolinkurl{inherited_ordinary_bridge.tex}.

\bibitem{salehi}
Saeed Salehi. Decidable formulas of intuitionistic primitive recursive
arithmetic. \emph{Reports on Mathematical Logic} \textbf{36} (2002), 55--61.
\url{https://rml.tcs.uj.edu.pl/rml-36/36-SALEHI.pdf}.
\end{thebibliography}
\end{document}
